
Predicting neural network accuracy from weights requires encoding weight tensors into fixed-size representations. Published comparisons between permutation-equivariant and flat MLP encoders have not controlled for encoder capacity, making it difficult to isolate the contribution of equivariant inductive bias from capacity differences. This work presents a controlled benchmark that matches three encoder types — flat MLP, Deep Sets (permutation-invariant), and Neural Functional Network (NFN, permutation-equivariant) — to approximately 500K parameters (±5\%) on the Schurholt MNIST-CNN model zoo (hyp_rand variant, 2,249 checkpoints). Downstream accuracy prediction quality is measured via Spearman rank correlation (ρ), and the capacity-wasting mechanism is measured directly via permutation sensitivity scores. The flat MLP achieves a sensitivity score of 0.649, while the NFN achieves 7.34 × 10⁻⁷ — a reduction of approximately 885,000×. On the test split (n = 338), the NFN achieves ρ = 0.6806 [95\% CI: 0.603, 0.748], the Deep Sets encoder achieves ρ = 0.4466 [95\% CI: 0.344, 0.544], and the untrained flat MLP baseline achieves ρ = 0.1688 [95\% CI: 0.069, 0.273], yielding Δρ(NFN − flat MLP) = 0.5119 [95\% CI: 0.381, 0.638]. A monotone ordering ρ(flat MLP) < ρ(Deep Sets) < ρ(NFN) is confirmed. Tier analysis reveals that NFN advantage is largest for low-accuracy models (ρ = 0.856) and inverts for high-accuracy models (ρ = −0.314), contrary to the pre-registered prediction of mid-tier dominance. The CIFAR-10 experiment was not executed due to a data download failure; all results are specific to the MNIST-CNN zoo. The flat MLP uses a single hidden layer of width 193 for a 2,464-dimensional input, which constitutes an architectural bottleneck that may inflate the observed Δρ.

